\section{Dominio applicativo: il segnale ECG}\label{sec:dominio-applicativo}

Il segnale elettrocardiografico (ECG) rappresenta la proiezione elettrica dell'attività cardiaca misurata sulla superficie corporea.
Dal punto di vista fisiologico, esso riflette la sequenza di depolarizzazione e ripolarizzazione delle camere cardiache e consente di osservare, con risoluzione temporale elevata, eventi che sono spesso immediatamente correlati allo stato elettrico del miocardio~\cite{malmivuo1995bioelectromagnetism}.

\subsection{Elettrofisiologia e morfologia del battito}

L'attivazione elettrica del cuore origina nel \textbf{nodo seno-atriale} (SA), che innesca la depolarizzazione degli atri.
Il fronte di onda raggiunge poi i ventricoli attraverso il \textbf{nodo atrioventricolare} (AV), il \textbf{fascio di His} e il sistema delle \textbf{fibre di Purkinje}, che distribuiscono l'impulso sul miocardio ventricolare in modo coordinato.
La sequenza temporale di queste fasi si traduce, sul tracciato, in segmenti e intervalli con significato clinico ben definito.

La morfologia tipica di un battito in ritmo sinusale è costituita dalle onde \textbf{P}, \textbf{QRS} e \textbf{T}:
l'onda P corrisponde alla depolarizzazione atriale; il complesso QRS alla depolarizzazione ventricolare (con i deflessi Q, R e S che descrivono l'attivazione sequenziale del setto e delle pareti ventricolari); l'onda T alla ripolarizzazione ventricolare.
Tra questi deflessi si misurano il \textbf{segmento PR} (conduzione atrioventricolare), il \textbf{segmento ST} (fase di ripolarizzazione precoce ventricolare) e l'\textbf{intervallo QT} (durata dell'attività ventricolare totale), spesso corretto per la frequenza cardiaca (QTc) nelle valutazioni cliniche.

La Figura~\ref{fig:ecg-morfologia} sintetizza, in forma schematica, la relazione tra fasi elettriche e morfologia osservabile.
In applicazioni automatiche, la forma locale del QRS e la durata degli intervalli costituiscono descrittori ad alta densità informativa per la distinzione tra battiti normali e patologici~\cite{dechazal2004automatic}.

% \begin{figure}[H]
%   \centering
%   \begin{tikzpicture}[
%     x=0.95cm, y=0.95cm,
%     >=Stealth,
%     wave/.style={line width=1.2pt, black},
%     guide/.style={dashed, gray!55, line width=0.4pt},
%     brace/.style={thick, black},
%     lbl/.style={font=\small},
%     dist/.style={font=\footnotesize}
%   ]
%     % --- assi ---
%     \draw[->] (-0.15,0) -- (11.5,0) node[below, font=\footnotesize] {tempo};
%     \draw[->] (0,-0.95) -- (0,1.55) node[left, font=\footnotesize, rotate=90, anchor=south] {ampiezza};
% 
%     % --- tracciato (coordinate fisse per allineare annotazioni) ---
%     \coordinate (Pon)  at (0.9,0);
%     \coordinate (Ppk)  at (1.35,0.24);
%     \coordinate (Poff) at (2.05,0);
%     \coordinate (Qon)  at (3.5,0);
%     \coordinate (Qpk)  at (3.65,-0.17);
%     \coordinate (Rpk)  at (4.15,1.08);
%     \coordinate (Spk)  at (4.55,-0.26);
%     \coordinate (Jpt)  at (4.95,0);    % fine QRS / inizio ST
%     \coordinate (Ton)  at (5.85,0);
%     \coordinate (Tpk)  at (7.0,0.36);
%     \coordinate (Toff) at (8.75,0);
% 
%     \draw[wave]
%       (0,0) -- (Pon)
%       .. controls (1.05,0.04) and (1.15,0.24) .. (Ppk)
%       .. controls (1.55,0.24) and (1.75,0.04) .. (Poff)
%       -- (Qon)
%       -- (Qpk) -- (Rpk) -- (Spk) -- (Jpt)
%       -- (Ton)
%       .. controls (6.2,0.06) and (6.6,0.38) .. (Tpk)
%       .. controls (7.4,0.38) and (8.2,0.06) .. (Toff)
%       -- (10.2,0);
% 
%     % --- etichette deflessioni (senza sovrapporre le frecce degli intervalli) ---
%     \node[lbl] at (Ppk)  [above=3pt] {P};
%     \node[lbl] at (Rpk)  [above=3pt] {R};
%     \node[lbl] at (Tpk)  [above=3pt] {T};
%     \node[lbl, anchor=east]  at ([xshift=-2pt]Qpk) {Q};
%     \node[lbl, anchor=west]  at ([xshift= 2pt]Spk) {S};
% 
%     % --- PR (sopra il tracciato) ---
%     \def\yPR{1.38}
%     \draw[guide] (Pon)  -- (Pon |- 0,\yPR);
%     \draw[guide] (Qon)  -- (Qon |- 0,\yPR);
%     \draw[brace, <->] (Pon |- 0,\yPR) -- node[above=4pt, dist] {PR} (Qon |- 0,\yPR);
% 
%     % --- durata QRS (sopra il complesso, tra inizio Q e punto J) ---
%     \def\yQRS{1.12}
%     \draw[guide] (Qon) -- (Qon |- 0,\yQRS);
%     \draw[guide] (Jpt) -- (Jpt |- 0,\yQRS);
%     \draw[brace, <->] (Qon |- 0,\yQRS) -- node[above=4pt, dist] {QRS} (Jpt |- 0,\yQRS);
% 
%     % --- segmento ST (tratto isoelettrico tra J e inizio T) ---
%     \def\yST{0.20}
%     \draw[guide] (Jpt |- 0,\yST) -- (Ton |- 0,\yST);
%     \node[dist, above=2pt] at ($(Jpt)!0.5!(Ton)$) {ST};
% 
%     % --- QT (sotto il tracciato, ben separato da Q e S) ---
%     \def\yQT{-0.72}
%     \draw[guide] (Qon)  -- (Qon |- 0,\yQT);
%     \draw[guide] (Toff) -- (Toff |- 0,\yQT);
%     \draw[brace, <->] (Qon |- 0,\yQT) -- node[below=4pt, dist] {QT} (Toff |- 0,\yQT);
%   \end{tikzpicture}
%   \caption{Morfologia schematica di un battito ECG in ritmo sinusale: deflessioni P, Q, R, S e T; intervalli PR e QT; durata del complesso QRS; segmento ST.}
%   \label{fig:ecg-morfologia}
% \end{figure}

Gli \emph{intervalli RR} (distanza temporale tra picchi R consecutivi) sono la misura principale per la frequenza cardiaca e per la \emph{Heart Rate Variability} (HRV).
In presenza di aritmie, la distribuzione degli RR perde regolarità: ectopie, pause o tachicardie producono sequenze che un classificatore può sfruttare insieme alla morfologia del battito isolato~\cite{taskforce1996hrv,dechazal2004automatic}.

\subsection{Variabilità della frequenza cardiaca}

La HRV quantifica le fluttuazioni dell'intervallo RR attorno al valore medio e riflette, in parte, l'equilibrio tra i sistemi simpatico e parasimpatico.
Le linee guida del Task Force del 1996~\cite{taskforce1996hrv} distinguono due famiglie di misure, entrambe rilevanti per il feature engineering adottato in questo lavoro.

Nel \textbf{dominio del tempo}, le metriche più diffuse includono:
\begin{itemize}
  \item \textbf{SDNN}: deviazione standard degli intervalli NN (battiti di origine sinusale);
  \item \textbf{RMSSD}: radice quadrata della media dei quadrati delle differenze successive tra intervalli NN;
  \item \textbf{pNN50}: percentuale di coppie di intervalli NN consecutivi la cui differenza supera 50~ms.
\end{itemize}
RMSSD e pNN50 sono particolarmente sensibili alle componenti ad alta frequenza, associate all'attività vagale a breve termine.

Nel \textbf{dominio della frequenza}, lo spettro della serie RR (tipicamente ricampionata e detrendata) viene suddiviso in bande:
la componente \textbf{LF} (circa 0.04--0.15~Hz) e la componente \textbf{HF} (circa 0.15--0.40~Hz), quest'ultima legata alla respirazione e al tono parasimpatico.
Il \textbf{rapporto LF/HF} è usato in letteratura come indicatore sintetico dell'equilibrio autonomico, sebbene la sua interpretazione clinica richieda cautela e contesto del paziente~\cite{taskforce1996hrv}.

Nel dataset impiegato, metriche HRV globali (SDNN, RMSSD, spettro LF/HF) non sono calcolate in tempo reale dalla pipeline Scala; il CSV contiene però intervalli RR locali (\texttt{pre-RR}, \texttt{post-RR}) e descrittori morfologici già estratti offline dal segnale grezzo.
Il classificatore opera quindi su un livello di astrazione che unifica morfologia istantanea e contesto ritmico immediato attorno al battito corrente.

\subsection{Rappresentazione computazionale e feature}

Dal punto di vista computazionale, il battito ECG può essere rappresentato come una \textbf{sequenza multiscala}: una componente temporale (distanza tra eventi) e una componente morfologica (forma del tracciato locale).
Questa doppia prospettiva consente di trasformare un segnale continuo in un insieme di osservazioni etichettate, ciascuna associata a un istante clinicamente significativo, adatto a modelli supervisionati che richiedono vettori numerici stabili nel tempo.

Nel file \texttt{MIT-BIH Arrhythmia Database.csv} del corpus Kaggle~\cite{sakib2021kaggle,sakib2021edge}, ogni riga descrive un singolo battito con \textbf{32 feature numeriche} (16 per ciascuna delle due derivazioni, prefissi \texttt{0\_} per lead~II e \texttt{1\_} per lead~V5), oltre a \texttt{record} (identificativo registrazione) e \texttt{type} (etichetta):
\begin{itemize}
  \item \textbf{Contesto ritmico}: \texttt{pre-RR} e \texttt{post-RR} (intervalli R--R rispettivamente precedente e successivo al battito corrente, in ms)~\cite{dechazal2004automatic}.
  \item \textbf{Morfologia delle onde}: ampiezze \texttt{pPeak}, \texttt{qPeak}, \texttt{rPeak}, \texttt{sPeak}, \texttt{tPeak}.
  \item \textbf{Durata degli intervalli}: \texttt{pq\_interval}, \texttt{qrs\_interval}, \texttt{st\_interval}, \texttt{qt\_interval}.
  \item \textbf{Forma del QRS}: cinque coefficienti \texttt{qrs\_morph0}--\texttt{qrs\_morph4} che riassumono la morfologia del complesso~\cite{clifford2006advanced}.
\end{itemize}

Tali attributi costituiscono il \emph{vettore di ingresso} per modelli di machine learning classico (Random Forest nel caso di studio) e sono prodotti da una catena di preprocessing che precede l'addestramento batch.
La scelta di non trasmettere il segnale raw riduce la complessità del modello, limita il volume dei messaggi Kafka e rende esplicita la semantica di ogni campo nel payload JSON.

\subsection{Preprocessing e qualità del segnale}\label{subsec:preprocessing}

La qualità del segnale condiziona in modo diretto l'affidabilità delle feature.
Prima della segmentazione beat-to-beat, la letteratura propone una sequenza di operazioni che, nel nostro flusso, sono applicate offline sulla fase di costruzione del CSV, ma che in un deployment edge replicherebbero lo stesso ordine logico~\cite{clifford2006advanced,pan1985realtime}.

\paragraph{Filtraggio.}
Un \textbf{filtro passa-alto} (tipicamente con frequenza di taglio intorno a 0.5~Hz) attenua il \emph{baseline wander}, derivante da respirazione, movimento del paziente o deriva degli elettrodi.
Un \textbf{filtro notch} centrato a 50~Hz o 60~Hz rimuove l'interferenza della rete elettrica (\emph{power-line interference}).
Un \textbf{filtro passa-basso} (spesso tra 35 e 40~Hz, coerente con la banda utile del MIT--BIH) limita il rumore muscolare e le componenti ad alta frequenza poco informative per la classificazione dei battiti.
In pratica, la combinazione di questi stadi equivale a un filtraggio passa-banda nell'intervallo 0.5--40~Hz citato nelle specifiche del database.

\paragraph{Rilevamento dei picchi R.}
L'algoritmo di \emph{Pan--Tompkins}~\cite{pan1985realtime} è un riferimento consolidato per la detection in tempo quasi reale.
La pipeline prevede: (i)~filtraggio passa-banda; (ii)~derivata per enfatizzare le pendenze del QRS; (iii)~elevamento al quadrato per uniformare la polarità; (iv)~integrazione a finestra mobile per smussare il segnale; (v)~soglia adattiva su due livelli per distinguere picchi veri da artefatti.
La robustezza dipende dalla calibrazione delle soglie e dalla qualità del tracciato: in presenza di aritmie con morfologie atipiche, un detector generico può richiedere post-processing o regole cliniche aggiuntive.

\paragraph{Indice di qualità del segnale.}
Per evitare che artefatti di movimento o contatto degradino le feature, si adottano \emph{Signal Quality Index} (SQI) che combinano metriche nel dominio del tempo, della frequenza o della forma d'onda~\cite{li2013sqi}.
Battiti associati a tratti di segnale con SQI insufficiente possono essere scartati, marcati come \emph{unknown} o inviati su un canale di quarantena, riducendo la propagazione di errori verso il classificatore.
In una pipeline distribuita, questa valutazione ha anche valenza architetturale: filtrare a monte significa trasmettere verso il broker solo eventi affidabili e alleggerire il carico del livello streaming.

Il preprocessing può dunque essere letto come \textbf{normalizzazione semantica}: il segnale fisiologico diventa una rappresentazione tabellare allineata al modello ML.
In produzione, la fase può risiedere sul dispositivo wearable o su un gateway, così da inviare al data center feature già validate anziché campioni raw ad alta frequenza.

\section{Dataset sperimentale: MIT--BIH tabellare (Kaggle)}\label{sec:dataset}

Per la fase sperimentale si utilizza la versione tabellare del \textbf{MIT--BIH Arrhythmia Database} pubblicata nel corpus Kaggle \emph{ECG Arrhythmia Classification Dataset}~\cite{sakib2021kaggle,sakib2021edge}, derivato dalla raccolta originale su PhysioNet~\cite{moody2001impact,goldberger2000physiobank}.
Il bundle Kaggle contiene più file CSV (MIT--BIH, MIT--BIH Supraventricular, INCART, SCD Holter); il prototipo Scala legge tramite \texttt{DATASET\_PATH} i file \texttt{*.csv} presenti nella directory configurata: nelle sperimentazioni descritte in questa tesi si assume il file \texttt{MIT-BIH Arrhythmia Database.csv}, così da restare allineati al benchmark MIT--BIH senza mescolare registrazioni di altri database.

\subsection{Origine e caratteristiche del corpus PhysioNet}

Il MIT--BIH originale~\cite{moody2001impact,goldberger2000physiobank} è una raccolta acquisita tra il 1975 e il 1979 presso il Beth Israel Hospital (oggi Beth Israel Deaconess Medical Center), con campionamento a 360~Hz e risoluzione di 11 bit su circa 10~mV di dinamica.
Il corpus è diventato, nel tempo, uno \emph{benchmark} de facto per la classificazione automatica dei battiti.

Le proprietà salienti del database grezzo sono le seguenti:
\begin{itemize}
  \item 48 registrazioni di circa 30 minuti ciascuna, da 47 soggetti (due tracce per un paziente).
  \item Due canali per registrazione, derivazioni differenti (tipicamente lead I e una derivazione del lead II modificato).
  \item Annotazioni beat-to-beat curate da più esperti, con codici simbolici per oltre venti categorie di battiti ed eventi non ritmici.
\end{itemize}

La versione CSV impiegata nel progetto non contiene campioni raw del segnale, bensì battiti già segmentati con feature estratte offline (filtraggio, rilevamento QRS, calcolo di intervalli e descrittori morfologici)~\cite{sakib2021edge,dechazal2004automatic}: la pipeline Scala legge questi archivi senza ripetere sul cluster l'intera catena descritta in \S\ref{subsec:preprocessing}.

\subsection{Schema del file CSV e volume dati}

Il file \texttt{MIT-BIH Arrhythmia Database.csv} contiene \textbf{100\,689} battiti da \textbf{44} registrazioni (colonna \texttt{record}), con \textbf{34} colonne totali (identificativo, etichetta e 32 feature).
Quattro tracce presenti nel MIT--BIH grezzo non compaiono nel CSV preprocessato, probabilmente escluse in fase di estrazione per qualità del segnale o annotazioni incomplete.
Ogni riga corrisponde a un battito; le feature duplicano la stessa semantica su due derivazioni (lead~II, prefisso \texttt{0\_}; lead~V5, prefisso \texttt{1\_}), come riassunto in Tabella~\ref{tab:csv-schema}.

\subsection{Schema di raggruppamento AAMI e etichette nel CSV}

Per confrontare i risultati con la letteratura clinica e con lo standard AAMI EC57:2012~\cite{aami2012ec57}, le annotazioni originali del MIT--BIH vengono raggruppate in cinque \textbf{super-classi}, come proposto da de~Chazal et al.~\cite{dechazal2004automatic} e adottato nel Capitolo~\ref{ch:analisi-dei-risultati-e-test}.
Nel CSV Kaggle, la colonna \texttt{type} adotta una nomenclatura leggermente diversa dalle sigle AAMI letterali, ma con corrispondenza biunivoca (Tabella~\ref{tab:aami-csv-mapping}).

\begin{table}[H]
  \centering
  \caption{Corrispondenza tra etichette nel CSV Kaggle e super-classi AAMI.}
  \label{tab:aami-csv-mapping}
  \begin{tabular}{@{}lll@{}}
    \toprule
    \textbf{CSV (\texttt{type})} & \textbf{AAMI} & \textbf{Significato} \\
    \midrule
    N    & N & Normale \\
    SVEB & S & Ectopia supraventricolare \\
    VEB  & V & Ectopia ventricolare \\
    F    & F & Fusione \\
    Q    & Q & Non classificabile \\
    \bottomrule
  \end{tabular}
\end{table}

\begin{table}[H]
  \centering
  \caption{Raggruppamento AAMI delle annotazioni MIT--BIH (sintesi del contenuto clinico).}
  \label{tab:aami-mapping}
  \begin{tabular}{@{}clp{7.5cm}@{}}
    \toprule
    \textbf{Classe} & \textbf{Sigla} & \textbf{Contenuto tipico} \\
    \midrule
    N & Normale & Battiti sinusali, blocchi di branca sinistro (L), blocchi di branca destro (R) \\
    S & Supraventricolare & Ectopie atriali (A), aberranti (a), nodali (J), ritmi supraventricolari \\
    V & Ventricolare & Contrazioni premature ventricolari (V), ectopie ventricolari (E) \\
    F & Fusion & Battiti di fusione (F) \\
    Q & Unknown & Non classificabili, rumore marcato (/), altre etichette non mappate \\
    \bottomrule
  \end{tabular}
\end{table}

La Tabella~\ref{tab:aami-mapping} non esaurisce ogni codice del database, ma fissa il vocabolario con cui si discutono metriche e errori clinici nel resto della tesi.
Lo \texttt{StringIndexer} di Spark ML indicizza i valori testuali effettivi (\texttt{N}, \texttt{SVEB}, \texttt{VEB}, \texttt{F}, \texttt{Q}) senza rinomina esplicita verso le sigle AAMI: la mappatura resta a carico dell'analisi dei risultati.

\begin{table}[H]
  \centering
  \caption{Feature per derivazione nel CSV \texttt{MIT-BIH Arrhythmia Database.csv} (prefisso \texttt{0\_} o \texttt{1\_}).}
  \label{tab:csv-schema}
  \begin{tabular}{@{}lp{8.5cm}@{}}
    \toprule
    \textbf{Suffisso colonna} & \textbf{Contenuto} \\
    \midrule
    \texttt{pre-RR}, \texttt{post-RR} & Intervalli R--R precedente e successivo (ms) \\
    \texttt{pPeak} \ldots \texttt{tPeak} & Ampiezze onde P, Q, R, S, T \\
    \texttt{pq\_interval}, \texttt{qrs\_interval}, \texttt{st\_interval}, \texttt{qt\_interval} & Durate intervalli/segm.\ (ms) \\
    \texttt{qrs\_morph0}--\texttt{qrs\_morph4} & Descrittori morfologici del complesso QRS \\
    \bottomrule
  \end{tabular}
\end{table}

\subsection{Sbilanciamento tra classi}

La distribuzione dei battiti nel file \texttt{MIT-BIH Arrhythmia Database.csv} è fortemente \emph{sbilanciata} (conteggi assoluti sul corpus completo):
N 90\,083 (89{,}5\%), VEB 7\,009 (7{,}0\%), SVEB 2\,779 (2{,}8\%), F 803 (0{,}8\%), Q 15 ($<$0{,}1\%).
Tale proporzione riflette la prevalenza clinica dei ritmi normali in registrazioni ambulatoriali selezionate, ma penalizza i modelli che ottimizzano l'accuratezza globale senza pesi per classe.

Durante l'addestramento si possono considerare \textbf{class weighting}, undersampling della classe maggioritaria o oversampling sintetico delle minoritarie; ciascuna strategia introduce trade-off (perdita di informazione, rischio di overfitting sulle classi rare, tempi di training più lunghi).
Nel prototipo attuale non è implementato alcun ribilanciamento esplicito: la scelta è coerente con l'obiettivo di dimostrare l'architettura end-to-end, ma va tenuta presente nell'interpretazione delle metriche del Capitolo~\ref{ch:analisi-dei-risultati-e-test}.

Dal punto di vista sperimentale, il MIT--BIH offre un compromesso favorevole tra complessità e riproducibilità: registrazioni sufficientemente eterogenee per stressare la pipeline, ma di dimensione gestibile in un contesto accademico in cui l'obiettivo principale non è il record di performance isolato, bensì la validazione di un flusso integrato dal dato storico allo streaming.

\subsection{Tipologie di anomalie contenute e limiti}\label{subsec:tipologie-di-anomalie}

Le annotazioni includono ectopie ventricolari (PVC), alterazioni di conduzione, battiti supraventricolari e, in alcune tracce, eventi rari o tratti rumorosi che mettono alla prova la segmentazione.
Il database non copre tutte le condizioni cliniche: aritmie prolungate come la fibrillazione atriale persistente o pattern legati a insufficienza cardiaca avanzata possono essere sottorappresentati.

La \emph{generalizzabilità} verso altre popolazioni (età, sesso, patologia, dispositivo di acquisizione) non va assunta per default.
I pazienti del MIT--BIH sono in gran parte adulti ricoverati o in follow-up cardiologico negli Stati Uniti della fine degli anni Settanta; un modello addestrato solo su questa distribuzione può degradare su cohort europee, pediatriche o su segnali da wearable consumer con banda e rumore differenti.

Per un classificatore supervisionato, la coerenza tra annotazioni, segmentazione e feature è tanto importante quanto l'algoritmo di apprendimento: errori a monte si propagano come etichette rumorose o descrittori distorti, limitando il soffitto di performance raggiungibile anche con modelli più complessi.

\section{Edge Computing e ruolo del simulatore}\label{sec:edge-computing}

Nel monitoraggio cardiaco continuo, la posizione dell'elaborazione lungo il percorso dato--cloud non è banale.
Si distinguono almeno tre livelli~\cite{kleppmann2017designing}:
\begin{itemize}
  \item \textbf{Cloud}: capacità computazionale elevata, storage centralizzato, latenza di rete non trascurabile;
  \item \textbf{Fog}: nodi intermedi (ospedale, edge data center) che aggregano flussi da più dispositivi;
  \item \textbf{Edge}: elaborazione sul wearable, sul Holter o sul gateway domestico, a ridosso del sensore.
\end{itemize}

Per l'ECG, spostare filtraggio, QRS detection e estrazione delle feature verso l'edge risponde a tre esigenze concrete: ridurre la latenza tra evento cardiaco e allarme; diminuire il traffico (un battito con 32 feature numeriche occupa ordini di grandezza in meno di un secondo di segnale a 360~Hz); limitare l'esposizione di dati biometrici raw oltre il perimetro controllato del paziente.

Dispositivi come Holter multicanale, patch ECG monouso e braccialetti con singola derivazione implementano già, in parte, questa logica: acquisiscono a alta frequenza localmente, poi trasmettono riassunti, eventi o tratti selezionati.
Il progetto di tesi non dispone di hardware dedicato; il componente \texttt{EcgPatientSimulator} emula il comportamento di un producer edge che pubblica feature già calcolate sul topic Kafka \texttt{ecg-input-stream}, in formato JSON.

\subsection{Funzione metodologica del simulatore}

Il simulatore consente di esercitare la pipeline end-to-end in laboratorio: stesso schema colonne del CSV di training, stesso topic, stesso consumer Spark, senza dipendenze da firmware o protocolli proprietari.
In uno scenario operativo, lo sostituirebbe un modulo on-device (o un gateway) con \texttt{KafkaProducer} o equivalente MQTT verso il broker.

L'approccio edge comporta vincoli noti: CPU e memoria limitate, consumo energetico, necessità di aggiornare modelli e soglie senza interrompere l'acquisizione.
Per la tesi, questi aspetti restano fuori dall'implementazione, ma orientano le scelte architetturali del Capitolo~\ref{ch:progettazione-e-architettura-del-sistema} (payload compatto, disaccoppiamento tramite Kafka, training batch separato dall'inferenza streaming).

La riproducibilità degli esperimenti ne beneficia: la sequenza di messaggi può essere rigenerata con la stessa logica (\texttt{limit(1000)}, intervallo di un secondo), permettendo confronti tra configurazioni del predictor o del checkpoint senza variabilità legata a dispositivi reali.

\subsection{Scenario evolutivo: digital twin}

A orizzonte più ampio, i flussi ECG da molteplici fonti potrebbero alimentare un \emph{digital twin} cardiaco del paziente: modello che integra storico clinico, eventi in tempo quasi reale e predizioni per simulare risposte a terapie o per prioritizzare alert.
Il prototipo descritto in questa tesi costituisce un mattone elementare di tale visione---ingestion, classificazione, persistenza versionata---senza pretendere di realizzarne la complessità clinica e regolatoria.

\section{Considerazioni etiche e implicazioni cliniche}\label{sec:considerazioni-etiche-e-implicazioni-cliniche}

Nell'applicazione di modelli ML a dati clinici, l'impatto di \textbf{falsi negativi} e \textbf{falsi positivi} va valutato prima delle metriche aggregate.
Un falso negativo su una ectopia ventricolare ripetuta può ritardare un intervento; un falso positivo aumenta il carico di revisione per il personale e, nel lungo periodo, riduce la fiducia nel sistema di alert.

Il problema è dunque operativo oltre che statistico: output, latenza e priorità degli allarmi devono essere allineati al contesto (telemetria domiciliare, reparto, terapia intensiva), dove i trade-off tra sensibilità e specificità non coincidono.

\paragraph{Quadro regolatorio.}
In Europa, software destinato a finalità mediche può rientrare nel campo del Regolamento (UE) 2017/745 sui dispositivi medici (MDR)~\cite{eu2017mdr}.
I sistemi di supporto alla diagnosi basati su algoritmi---spesso indicati come \emph{Software as a Medical Device} (SaMD)---richiedono gestione del rischio, tracciabilità, validazione clinica e monitoraggio post-market.
Il prototipo di questa tesi non è certificato né destinato all'uso clinico; le considerazioni regolatorie servono a inquadrare il divario tra dimostrazione accademica e prodotto deployabile.

\paragraph{Validazione, interpretabilità e bias.}
Un'adozione responsabile presuppone validazione su popolazioni e dispositivi diversi da quelli del training, documentazione delle prestazioni per sottogruppo e strumenti di interpretabilità (importanza delle feature, analisi degli errori per classe AAMI).
Il MIT--BIH, per quanto prezioso, introduce \emph{bias} demografico e tecnologico: campione storico, non necessariamente rappresentativo dell'utenza attuale né delle acquisizioni da wearable.

Infine, la pipeline deve essere leggibile da chi risponde dei dati: finalità del trattamento, minimizzazione (feature anziché raw dove possibile), retention su Delta Lake e accessi al broker vanno progettati con pari attenzione rispetto all'accuratezza del classificatore.
Senza questi accorgimenti, anche un modello accurato resterebbe confinato al laboratorio.
